%! Sine, Cosine, and Tangent — Unit 3, Lesson 3.2 — Homework (Answer Key)
\documentclass[10pt]{article}
\usepackage{precalculus-article}
\usepackage{precalculus-key}   % pulls in -boxes; do NOT also load -boxes
\newcommand{\TallMath}[1]{$\displaystyle #1\rule[-1.4em]{0pt}{3.2em}$}

\begin{document}

\pageheader{Unit 3, Lesson 3.2}{Homework --- Answer Key}
\namedateperiod

\vspace{0.10in}

\begin{notesbox}{Homework: Sine, Cosine, and Tangent --- Key}

\textbf{1.} Exact values table:

\medskip
\begin{center}
\renewcommand{\arraystretch}{2.4}
\begin{tabular}{|c|c|c|c|c|}
\hline
$\theta$ & Degrees & $\sin\theta$ & $\cos\theta$ & $\tan\theta$ \\
\hline
$0$     & $0°$  & \ans{$0$} & \ans{$1$} & \ans{$0$} \\
$\pi/6$ & $30°$ & \ans{$\dfrac{1}{2}$} & \ans{$\dfrac{\sqrt{3}}{2}$} & \ans{$\dfrac{\sqrt{3}}{3}$} \\
$\pi/4$ & $45°$ & \ans{$\dfrac{\sqrt{2}}{2}$} & \ans{$\dfrac{\sqrt{2}}{2}$} & \ans{$1$} \\
$\pi/3$ & $60°$ & \ans{$\dfrac{\sqrt{3}}{2}$} & \ans{$\dfrac{1}{2}$} & \ans{$\sqrt{3}$} \\
$\pi/2$ & $90°$ & \ans{$1$} & \ans{$0$} & \ans{undefined} \\
\hline
\end{tabular}
\end{center}

\vspace{0.14in}

\textbf{2.} $P = (-0.28,\; 0.96)$:

\ans{$\sin\theta = 0.96$, \quad $\cos\theta = -0.28$, \quad
$\tan\theta = 0.96/(-0.28) \approx -3.43$}

\vspace{0.12in}

\textbf{3.} Why is $\tan(\pi/2)$ undefined?

\ans{At $\theta = \pi/2$, the terminal ray points straight up and $P = (0, 1)$, so $\cos(\pi/2) = 0$.
Since $\tan\theta = \sin\theta / \cos\theta$, and division by zero is undefined, $\tan(\pi/2)$ is
undefined.}

\vspace{0.12in}

\textbf{4.} AP-style MC answer: \ans{(B) $-3/5$}

\ans{In Quadrant III, both $x$ and $y$ are negative, so $\cos\theta < 0$.
Only option (B) is negative and has absolute value less than 1 (required since $|\cos\theta| \le 1$).}

\vspace{0.12in}

\textbf{5.} Given $\cos\theta = -4/5$ and $\sin\theta > 0$, find $\tan\theta$:

\ans{Step 1: Use $\sin^2\theta + \cos^2\theta = 1$.\\[3pt]
$\sin^2\theta + (-4/5)^2 = 1 \;\Rightarrow\; \sin^2\theta = 1 - 16/25 = 9/25$.\\[3pt]
Step 2: Since $\sin\theta > 0$, take the positive root: $\sin\theta = 3/5$.\\[3pt]
Step 3: $\tan\theta = \sin\theta / \cos\theta = (3/5)/(-4/5) = -3/4$.}

\begin{teachernote}
Problem 5 is a critical skill for AP Precalculus. The key error students make is forgetting to
use the given sign condition ($\sin\theta > 0$) when taking the square root in Step 2.

After computing $\sin^2\theta = 9/25$, prompt students: ``$\sin\theta = \pm 3/5$ --- which sign?''
Then direct them to the given condition. This mirrors what students will encounter on AP-style
free-response items.

Note: $\cos\theta = -4/5$ with $\sin\theta > 0$ places $\theta$ in Quadrant II, where $x < 0$ and
$y > 0$. You can ask students to confirm this makes sense with $\tan\theta = -3/4 < 0$ (tangent is
negative in Quadrant II).
\end{teachernote}

\end{notesbox}

\begin{extensionbox}
\textbf{Extension: Deriving the Pythagorean Identity for Tangent --- Key}

\begin{enumerate}[leftmargin=1.8em, itemsep=4pt]
  \item Dividing $\sin^2\theta + \cos^2\theta = 1$ by $\cos^2\theta$:
        \[
          \frac{\sin^2\theta}{\cos^2\theta} + 1 = \frac{1}{\cos^2\theta}
        \]
  \item \ans{$\dfrac{\sin^2\theta}{\cos^2\theta} = \left(\dfrac{\sin\theta}{\cos\theta}\right)^2 = \tan^2\theta$}
  \item \ans{The identity is $\tan^2\theta + 1 = \sec^2\theta$, where $\sec\theta = 1/\cos\theta$.}
\end{enumerate}
\end{extensionbox}

\vspace{0.10in}

\begin{notesbox}{Preview: Lesson 3.3}
In \textbf{Lesson 3.3} we derive the exact coordinates of $P$ for
$\theta = \pi/6,\, \pi/4,\, \pi/3$ using 45-45-90 and 30-60-90 right-triangle ratios.
Your completed table in Problem 1 will be fully confirmed with geometric proofs.
\end{notesbox}

\end{document}
